\begin{tikzpicture}[x=6.5mm, y=6.5mm]
  % trzy siatki 5 x 5; komórka (i, j) wyróżniona, gdy warunek jest prawdziwy
  \foreach \g/\x/\warunek in {0/0/{\j>=\i}, 1/7.5/{(\i==0)||(\i==4)||(\j==0)||(\j==4)}, 2/15/{mod(\i+\j,2)==0}} {
    \foreach \j in {0,...,4}
      \node[indeks] at (\x+\j, 0.75) {\j};
    \foreach \i in {0,...,4} {
      \node[indeks] at (\x-0.85, -\i) {\i};
      \foreach \j in {0,...,4} {
        \pgfmathtruncatemacro{\on}{(\warunek) ? 1 : 0}
        \ifnum\on=1
          \node[komorka wyr, minimum width=6.5mm, minimum height=6.5mm] at (\x+\j, -\i) {*};
        \else
          \node[komorka, draw=linia, minimum width=6.5mm, minimum height=6.5mm] at (\x+\j, -\i) {};
        \fi
      }
    }
    \node[indeks, text=akcent] at (\x+2, 1.35) {kolumna \kod{j}};
    \node[indeks, text=akcent, rotate=90] at (\x-1.45, -2) {wiersz \kod{i}};
  }
  \node[kod, font=\ttfamily\footnotesize, align=center, anchor=north] at (2, -5.05) {j >= i};
  \node[kod, font=\ttfamily\footnotesize, align=center, anchor=north] at (9.5, -5.05) {i == 0 or i == n - 1\\or j == 0 or j == n - 1};
  \node[kod, font=\ttfamily\footnotesize, align=center, anchor=north] at (17, -5.05) {(i + j) \% 2 == 0};
  \node[podpis] at (2, -6.6) {górny trójkąt};
  \node[podpis] at (9.5, -6.6) {ramka};
  \node[podpis] at (17, -6.6) {szachownica};
\end{tikzpicture}
